% =========================================================================
\section{Experiments}
\label{sec:experiments}

% The experiments follow three regimes of increasing observation
% difficulty (same task, same model throughout):
%   Regime I   — isolated plant, single clean view
%   Regime II  — morphological complexity, viewpoint, self-occlusion
%   Regime III — neighbor occlusion in realistic rows
% This progression is the narrative; the task formulation (any subset ->
% procedural representation) is unchanged.

\subsection{Setup}
\label{sec:exp-setup}
% - PER-TRAIT metrics, never one pooled MAE: leaf count exact / ±1
%   accuracy; height and leaf length/width relative error; branching and
%   insertion angles MAE in degrees; azimuth circular MAE; droop/curl/
%   twist normalized angular error. Every trait compared against three
%   trivial baselines: dataset mean, plant mean, position-only predictor;
%   traits that fail to beat them are MARKED, not averaged in.
% - Geometry: Chamfer AND density-aware CD \citep{wu2021dcd} AND
%   F-score@tau (slender-aware); coverage@threshold.
% - Representation: frozen-probe R^2 (ridge on frozen CLS, train-fit
%   standardizer).
% - Baselines: completion (AdaPoinTr \citep{yu2023adapointr}, EINet
%   \citep{cai2024einet}) + our param head on their output (two-stage);
%   Demeter/Ghrer \citep{demeter2025,ghrer2025paramplants} where
%   adaptable; staged VLM pipeline \citep{hadadi2026vlmmaize} (isolated
%   plants only); frozen encoders + supervised-from-scratch.
% - Species roles (per generator statistics): MAIZE carries the
%   independent-trait recovery claims; SORGHUM is cross-species
%   replication (its height/leaf-count/biomass are one generative
%   variable and are never presented as three findings).
% - Implementation: \todo{114M encoder; 197,400-step budget}.
\todo{compact setup prose; implementation details to supplement}

\subsection{Regime I: single-view reconstruction of isolated plants}
\label{sec:exp-regime1}
% The clean headline: can one RGB observation of an isolated plant
% constrain an editable procedural model? Evidence in hand: zero-shot
% maize generation from one photograph at ~1.5x the masked-
% reconstruction error, before any distillation.
\begin{table}[h]\centering\small
\begin{tabular}{lccccc}
\toprule
Input & Leaf cnt & Height & Leaf params & F-score & Chamfer \\
 & (acc$\pm$1) & (rel.\ err) & (mean skill) & $\uparrow$ & $\downarrow$ \\
\midrule
RGB (1 view)      & \todo{} & \todo{} & \todo{} & \todo{} & \todo{} \\
RGB (3 views)     & \todo{} & \todo{} & \todo{} & \todo{} & \todo{} \\
Depth (1 view)    & \todo{} & \todo{} & \todo{} & \todo{} & \todo{} \\
PC (1 view)       & \todo{} & \todo{} & \todo{} & \todo{} & \todo{} \\
RGB+D+PC          & \todo{} & \todo{} & \todo{} & \todo{} & \todo{} \\
\midrule
Two-stage (completion$\to$fit) & \todo{} & \todo{} & \todo{} &
\todo{} & \todo{} \\
\bottomrule
\end{tabular}
\caption{\textbf{Regime I (isolated plants, maize).} Single-view inputs
against the multi-view and multimodal upper bounds; one amortized model
against two-stage completion-then-fit. Distilled students throughout;
before/after distillation appears in Tab.~\ref{tab:main}.
\todo{fill; sorghum replication to supplement}}
\label{tab:regime1}
\end{table}
% ---- FIG 3 (qualitative, full width): input photograph -> predicted
% parameters -> re-rendered plant vs GT, including ONE EDIT demo (change
% leaf angle, re-render) proving the output is genuinely procedural.
\todo{Fig 3 qualitative + edit demo + 2-3 sentences of interpretation}

\subsection{Regime II: morphological complexity, viewpoint,
self-occlusion}
\label{sec:exp-regime2}
% - Complexity sweep: leaf count 3 -> 15(+) on the x-axis (generator
%   already spans 3-22); curves per view configuration; y = per-trait
%   error / F-score. THE question: when does single-view procedural
%   inversion become ambiguous?
% - Camera study with EVIDENCE-BASED axes (not generic 0/45/90 azimuth):
%   (a) ELEVATION ladder — sorghum shows difficulty is set by elevation
%   (U-curve, best side-on); (b) azimuth RELATIVE TO THE LEAF PLANE for
%   distichous maize (the axis the pose forensics exposed). Multi-view
%   fusion as approximate upper bound; a small camera schematic figure.
% - View-consistency metric: ten predictions of one plant's parameters;
%   report normalized across-view std alongside accuracy. State
%   explicitly that STABLE IS NOT RIGHT (zero-skill fields are stable
%   at the dataset mean; distillation makes skilled fields both more
%   accurate and more stable, App.~\ref{app:ablations}).
% ---- FIG 4: complexity sweep (x = leaf count, curves = view configs).
\todo{Fig 4 + camera schematic + prose}

\subsection{Regime III: neighbor occlusion in realistic rows}
\label{sec:exp-occlusion}
% - PAIRED design: every evaluation plant rendered twice — isolated, and
%   inserted into a row with neighbors — identical morphology, only
%   observation difficulty changes. Report the occlusion-induced
%   increase in each per-trait error directly.
% - Row spacing swept continuously (\todo{0.15-0.70 m}); occlusion
%   severity REPORTED CONTINUOUSLY, not as a binary: visible pixel
%   fraction, visible 3D-point fraction, neighbor-point contamination
%   inside the crop region, fraction of leaves partially/fully hidden.
% - FIG 5: per-trait error and F-score vs visible fraction; ours vs
%   AdaPoinTr/EINet two-stage (completion baselines hallucinate neighbor
%   geometry into the target; procedural output degrades gracefully).
% - Which traits survive occlusion, which become ambiguous; per-leaf
%   analysis (lowest, most-occluded leaves fail first).
\todo{Fig 5 + paired-occlusion table + prose}

\subsection{Any-subset inference and distillation}
\label{sec:exp-main}
% Full subset grid + the distillation axis (kept in the main paper: the
% +19.1%/-43.8% numbers carry the deployment claim).
\begin{table}[h]\centering\small
\begin{tabular}{lcccc}
\toprule
Input & \multicolumn{2}{c}{Param err $\downarrow$} &
\multicolumn{2}{c}{Geometry $\downarrow$} \\
 & before & after distill & before & after \\
\midrule
RGB only        & \todo{} & \todo{} & \todo{} & \todo{} \\
Depth only      & \todo{} & \todo{} & \todo{} & \todo{} \\
PC only         & \oldnum{0.062} & \oldnum{0.035} & \todo{} & \todo{} \\
RGB+D           & \todo{} & \todo{} & \todo{} & \todo{} \\
RGB+PC          & \todo{} & \todo{} & \todo{} & \todo{} \\
All three       & \todo{} & \todo{} & \todo{} & \todo{} \\
\bottomrule
\end{tabular}
\caption{\textbf{Any-subset results, before and after distillation from
the vision-complete teacher} (RGB+D+PC; the teacher never observes
parameters, Sec.~\ref{sec:method-distill}). Distillation concentrates
its gains on the weakest sources. \todo{fill; both species}}
\label{tab:main}
\end{table}

\subsection{Where the parameter stream belongs}
\label{sec:exp-params}
% Precise claim (not the slogan): procedural parameters are valuable
% prediction targets, but treating them as an auxiliary representation-
% learning modality does not improve, and can degrade, the sensor
% representation.
% - E2 table (4 arms) with probe R^2 columns BESIDE chamfer.
% - One sentence each: token-budget control arms (budget restored, gap
%   persists, dose-dependent); parameters-as-target-only (still below
%   PC+RGB+D).
% - The parameter-INPUT teacher appears here as the privileged-
%   information ablation (Sec.~\ref{sec:method-distill}).
% - Mechanism sentence, hedged \citep{peng2022ogm,yu2020pcgrad}.
% Full grid -> App.~\ref{app:ablations}.
\todo{Table (E2 + probe) + prose}

\subsection{What the metrics disagree about}
\label{sec:exp-metrics}
% - Modality ranking by Chamfer vs by probe (depth: -32% vs ~0).
% - FIG 6: a reconstruction with good Chamfer that misses a blade, vs one
%   with worse Chamfer that recovers structure (the slender-blindness
%   exhibit); plus scaling curves (chamfer keeps paying, probe saturates
%   at ~3K plants).
% - Decision rule: parameter recovery is PRIMARY for the inverse-
%   procedural task; geometry metrics are complementary (re-render
%   fidelity). Probe decides modality questions.
\todo{Fig 6 + prose}

\subsection{Real-world validation}
\label{sec:exp-real}
% - PRIMARY: real field scans. Isolated real plants: ours (each input
%   subset) vs the staged VLM pipeline \citep{hadadi2026vlmmaize}
%   (accuracy parity question + speed/cost: pipeline ~600 s + \$1.30 per
%   plant; ours \todo{ms}). Occluded real rows: ours only — the pipeline
%   cannot run there. Zero-shot and fine-tuned rows.
% - SUPPLEMENT: targeted greenhouse study with matched acquisition
%   (strobe/short-exposure imaging if adopted — see Sec.~
%   \ref{sec:method-data} background note), manually measured traits
%   (leaf count, height, representative lengths/widths, insertion
%   angles), a small high-quality-scan subset. Synthetic-only training
%   -> real RGB -> recovered traits.
% - TABLE: accuracy + seconds/plant + $/plant. FIG 7: real examples
%   including an occluded row.
\todo{Table + Fig 7 + prose}

\subsection{Representation quality and parsimony}
\label{sec:exp-repr}
% - Frozen-encoder table (ours vs EmbodiedMAE-B/Point-MAE/DINOv2/
%   MultiMAE/supervised) — fusion, not a better PC encoder, is the edge;
%   camera-pose forensics one-liner (gravity-only: 5/9).
% - Compression: bytes/plant procedural vs Draco/G-PCC/3DAC
%   \citep{fang2022threedac} at matched geometry + what codecs CANNOT do
%   (edit, simulate, read traits). TALKING POINT: ~3 orders of magnitude.
\todo{Table 5 + Table 6 + prose}
